% Timing diagram of a three-instruction polling loop going round once.
%
%       loop  MOVE  @R0, R2        ; 0x0005          (1 word)
%             AND   0x0002, R2     ; 0x0006, 0x0007  (2 words)
%             RBRA  loop, Z        ; 0x0008, 0x0009  (2 words)
%
% Five words of instruction memory, and exactly ten clock cycles per iteration.
% The loop waits for bit 1 of an I/O device mapped at the address in R0; here
% that device is an ordinary memory word (0x000A) holding zero, so the branch is
% always taken and the loop runs forever. See doc/README.md for the prose.
%
% All values below are taken from a GHDL simulation of test/prog_poll.asm,
% sampled at cycles 35..45 of the run (350 ns .. 450 ns), by which point the
% loop has settled: cycle 45 is bit-for-bit identical to cycle 35, which is why
% the last column of the diagram repeats the first.
%
% Drawing convention, the same one src/cpu_main/timing.tex uses: handshake bits
% (valid/ready/enable) are always drawn with their true value; payload buses are
% greyed out whenever their valid is low, and greyed out but still labelled when
% they carry a word of the abandoned, speculatively fetched stream past the
% branch (0x000A and up). What is added here is COLOUR: every box is tinted
% according to which of the three instructions it belongs to, so that the ten
% cycles can be attributed to them by eye.
%
% The top three rows are not signals but a summary of the three pipeline
% registers below them, each derived from the valid/address pair drawn further
% down: ICACHE/m_valid_o + m_addr_o, DECODE/seq_valid_o + seq_stage_o.addr,
% and WRITE/prep_valid_i + prep_stage_i.addr. A low line means the stage is
% empty in that cycle.
%
% WRITE/prep_stage_i.microcode is the single 12-bit micro-operation SEQUENCER
% selected out of DECODE's list, i.e. the one being executed. 0x084 is
% MEM_READ_SRC + REG_MOD_SRC, 0x828 is LAST + MEM_WAIT_SRC + REG_WRITE, and
% 0x808 is LAST + REG_WRITE.

\documentclass{standalone}
\usepackage{timing}

% Eleven cycles of four-character bus values, plus the widest label in the tree.
\renewcommand{\wavescale}{1.4}

% One tint per instruction, so a box says at a glance which of the three owns
% the cycle. Deliberately pale: the value inside is printed in black.
\newcommand*{\cmove}[2]{\bitvector{#1}{#2}{blue!15}}
\newcommand*{\cand}[2]{\bitvector{#1}{#2}{green!25}}
\newcommand*{\crbra}[2]{\bitvector{#1}{#2}{orange!35}}

\begin{document}
\begin{wave}{29}{11}
 \nextwave{\textit{in ICACHE}}              \bit{0}{3} \cmove{MOVE}{1} \cand{AND}{3} \crbra{RBRA}{2} \unknown[x]{1} \bit{0}{1}
 \nextwave{\textit{in DECODE}}              \unknown[x]{1} \bit{0}{3} \cmove{MOVE}{3} \cand{AND}{1} \bit{0}{1} \crbra{RBRA}{1} \unknown[x]{1}
 \nextwave{\textit{in WRITE}}               \crbra{RBRA}{1} \bit{0}{4} \cmove{MOVE}{1} \bit{0}{1} \cmove{MOVE}{1} \cand{AND}{1} \bit{0}{1} \crbra{RBRA}{1}
 \nextwave{FETCH/wb\_stb\_o}                \bit{1}{6} \bit{0}{1} \bit{1}{4}
 \nextwave{FETCH/wb\_addr\_o}               \unknown[000D]{1} \cmove{0005}{1} \cand{0006}{1} \cand{0007}{1} \crbra{0008}{1} \crbra{0009}{1} \unknown[000A]{2} \unknown[000B]{1} \unknown[000C]{1} \unknown[000D]{1}
 \nextwave{FETCH/dc\_valid\_o}              \bit{0}{2} \bit{1}{8} \bit{0}{1}
 \nextwave{FETCH/dc\_addr\_o}               \unknown[x]{2} \cmove{0005}{1} \cand{0006}{1} \cand{0007}{1} \crbra{0008}{2} \crbra{0009}{1} \unknown[000A]{1} \unknown[000B]{1} \unknown[x]{1}
 \nextwave{ICACHE/s\_ready\_o}              \bit{0}{1} \bit{1}{4} \bit{0}{1} \bit{1}{4} \bit{0}{1}
 \nextwave{ICACHE/m\_valid\_o}              \bit{0}{3} \bit{1}{7} \bit{0}{1}
 \nextwave{ICACHE/m\_ready\_i}              \bit{1}{4} \bit{0}{2} \bit{1}{1} \bit{0}{1} \bit{1}{3}
 \nextwave{ICACHE/m\_double\_o}             \bit{0}{5} \bit{1}{2} \bit{0}{1} \bit{1}{1} \bit{0}{2}
 \nextwave{ICACHE/m\_addr\_o}               \unknown[x]{3} \cmove{0005}{1} \cand{0006}{3} \crbra{0008}{2} \unknown[000A]{1} \unknown[x]{1}
 \nextwave{DECODE/seq\_valid\_o}            \bit{1}{1} \bit{0}{3} \bit{1}{4} \bit{0}{1} \bit{1}{2}
 \nextwave{DECODE/seq\_ready\_i}            \bit{1}{4} \bit{0}{2} \bit{1}{5}
 \nextwave{DECODE/seq\_stage\_o.addr}       \unknown[000A]{1} \unknown[x]{3} \cmove{0005}{3} \cand{0006}{1} \unknown[x]{1} \crbra{0008}{1} \unknown[000A]{1}
 \nextwave{WRITE/prep\_valid\_i}            \bit{1}{1} \bit{0}{4} \bit{1}{1} \bit{0}{1} \bit{1}{2} \bit{0}{1} \bit{1}{1}
 \nextwave{WRITE/prep\_stage\_i.addr}       \crbra{0008}{1} \unknown[x]{4} \cmove{0005}{1} \unknown[x]{1} \cmove{0005}{1} \cand{0006}{1} \unknown[x]{1} \crbra{0008}{1}
 \nextwave{WRITE/mem\_req\_valid\_o}        \bit{0}{5} \bit{1}{1} \bit{0}{5}
 \nextwave{MEMORY/msrc\_valid\_o}           \bit{0}{6} \bit{1}{1} \bit{0}{4}
 \nextwave{REGISTERS/wr\_en\_i}             \bit{1}{1} \bit{0}{6} \bit{1}{2} \bit{0}{1} \bit{1}{1}
 \nextwave{REGISTERS/wr\_reg\_i}            \crbra{F}{1} \unknown[x]{6} \cmove{2}{1} \cand{2}{1} \unknown[x]{1} \crbra{F}{1}
 \nextwave{CPU/inst\_done\_o}               \bit{1}{1} \bit{0}{6} \bit{1}{2} \bit{0}{1} \bit{1}{1}
 \nextwave{WRITE/fetch\_valid\_o}           \bit{1}{1} \bit{0}{9} \bit{1}{1}
 \nextwave{WRITE/fetch\_addr\_o}            \crbra{0005}{1} \unknown[x]{9} \crbra{0005}{1}
\end{wave}

\end{document}
